% !TEX program = xelatex
\documentclass{article}
\usepackage{graphicx} % Required for inserting images
\usepackage{amsmath, amsthm, amssymb, amsfonts}
\usepackage{geometry}
\usepackage{array}
\usepackage{tikz}
\usepackage{hyperref}
\hypersetup{colorlinks=true, linkcolor=blue}
\usepackage{xcolor}
\usepackage{longtable}
\usepackage{fontspec}
\theoremstyle{plain}
\usepackage{amsthm}
% Style for core mathematical results
\theoremstyle{plain}
\newtheorem{thm}{Theorem}[section]
\newtheorem{lem}[thm]{Lemma}
\newtheorem{propn}[thm]{Proposition}
\newtheorem{cor}[thm]{Corollary}
% Style for formal definitions and contextual items
\theoremstyle{definition}
\newtheorem{defn}{Definition}[section]
\newtheorem{eg}{Example}[section]
\newtheorem{conj}{Conjecture}[section]
\newtheorem{prblm}{Problem}[section]
% Style for informal or supplementary notes
\theoremstyle{boldremark}
\newtheorem*{rk}{Remark}
\newtheorem*{note}{Note}

\usepackage{eso-pic, xcolor, graphicx}


\newfontfamily{\tamilfont}{Nirmala UI}
\usetikzlibrary{arrows.meta, decorations.pathmorphing, calc, positioning}
% 1. Link the 'nickname' \customzha to your specific file
\newfontfamily{\customzha}[Path = ./]{HindMadurai-Light.ttf}
% 2. Use the nickname to change the font for just that character
\renewcommand{\qedsymbol}{\textcolor{red}{{\customzha ழ}}}
\DeclareMathOperator{\im}{im}
\definecolor{cyan}{HTML}{00FFFF}



\geometry{a4paper, margin=1in}

\title{\textbf{Stochastic Differential Equations}\thanks{This is a preliminary draft and currently a work in progress.}}
\author{Gowriprakash}

\begin{document}

\maketitle

15/05/26\\
the subject seem to be using a lot of constructs from measure theory and probability theory, the intersection being the probability measure.

16/05/26\\
i'm reading the papers by song and karras.

17/05/26\\

20/05/26\\
i was getting a bit mathsy with rvs and dr. p clearly said to be keep it physicy and the next week's abjective is to understand the fokker planck equation. i have started with the book by prof. v. balakrishnan on elements of non equilibrium statistical mechanics. 

26/05/26\\
i think i've reached a point where i can start writing the exposition.

\section*{Langevin equation}

The objective of this equation is to model the motion of a particle ("tagged particle" as we call it) in a fluid, which is an elementary problem in the branch of non equilibrium statistical mechanics. \\
\\
The equation can be systematically obtained from the newton's 2nd law as follows. Consider a tagged particle of mass $m$ in a fluid which is being acted upon by an external force $\mathbf{F}$. Now newton's 2nd law reads,
\begin{equation}
    m \dot{\mathbf{v}} = \mathbf{F} 
\end{equation}
Now we split our system into 2 parts, the tagged particle and the rest of the system (i.e., the fluid). 



%references:
%1. Song, J., Sohl-Dickstein, J., Kingma, D. P., Kumar, A., Ermon, S., & Poole, B. (2021). Score-based generative modeling through stochastic differential equations. arXiv preprint arXiv:2011.13456.
%2. Karras, T., Laine, S., Aittala, M., Hellsten, J., Lehtinen, J., & Aila, T. (2022). Elucidating the design space of diffusion-based generative models. arXiv preprint arXiv:2206.00364.
%3. stochastic differential equations: an introduction with applications by bernt oksendal
%4. an intro to real and complex manifolds, giuliano sorani
\end{document}